\par\Needspace{6\baselineskip}\medskip
\begin{center}\hypertarget{native-226-article163}{}\pdfbookmark[2]{163.}{article-163}163.\end{center}
The method we have set out since Article 136 is adapted chiefly to the first determination of a hitherto entirely unknown orbit. It can also be used with equal success, however, when an already approximately known orbit is to be corrected using three observations however far apart. In such a case it is advisable to make certain changes. Whenever the observations cover a very large heliocentric motion, we can no longer take $\dfrac{\theta''}{\theta}$ and $\theta\theta''$ as approximate values of the quantities $P$, $Q$. Much more accurate values can instead be derived from the approximately known elements. From those elements we can thus readily calculate the heliocentric positions in the orbit for the times of the three observations. If we denote the true anomalies by $v$, $v'$, $v''$, the radius vectors by $r$, $r'$, $r''$, and the semiparameter by \SourceLeaf{227}$p$, we obtain the following approximate values.
\[
\begin{aligned}
P&=\frac{r\sin(v'-v)}{r''\sin(v''-v')}\mathpunct{\raisebox{.45ex}{,}}\\[.6em]
Q&=\frac{4r'^4\sin\frac12(v'-v)\sin\frac12(v''-v')}{p\cos\frac12(v''-v)}\mathpunct{\cdot}
\end{aligned}
\]
\noindent We thus base the first hypothesis on these values, and the second and third on values slightly varied at will. It would be inappropriate here to adopt $P'$ and $Q'$ as new values (as we did above), since we can no longer assume that these values will be more accurate. In this way all three hypotheses can most conveniently be worked out \emph{at once}. The fourth will then be formed according to the directions in Article 120. We would have no objection, however, if someone found one or another of the ten methods set out in Articles 124---129 at least nearly as convenient in such a case, if not more so, and therefore preferred to use it.
\par\bigskip\begin{center}\rule{.23\linewidth}{.35pt}\end{center}
